% Credits are indicated where needed. The general idea is based on a template by Vel (vel@LaTeXTemplates.com) and Frits Wenneker.
%\documentclass[UTF8]{ctexart}
\documentclass[11pt, letterpaper]{article} % General settings in the beginning (defines the document class of your paper)
% 11pt = is the font size
% A4 is the paper size
% “article” is your document class

%----------------------------------------------------------------------------------------
%	Packages
%----------------------------------------------------------------------------------------
\usepackage[UTF8]{ctex}
% Necessary
\usepackage[german,english]{babel} % English and German language
\usepackage{booktabs} % Horizontal rules in tables
% For generating tables, use “LaTeX” online generator (https://www.tablesgenerator.com)
\usepackage{comment} % Necessary to comment several paragraphs at once
%\usepackage[utf8]{inputenc} % Required for international characters
\usepackage[T1]{fontenc} % Required for output font encoding for international characters
\usepackage{listings}
% Might be helpful
\usepackage{amsmath,amsfonts,amsthm} % Math packages which might be useful for equations
\usepackage{tikz} % For tikz figures (to draw arrow diagrams, see a guide how to use them)
\usepackage{tikz-cd}
\usetikzlibrary{positioning,arrows} % Adding libraries for arrows
\usetikzlibrary{decorations.pathreplacing} % Adding libraries for decorations and paths
\usepackage{tikzsymbols} % For amazing symbols ;) https://mirror.hmc.edu/ctan/graphics/pgf/contrib/tikzsymbols/tikzsymbols.pdf
\usepackage{blindtext} % To add some blind text in your paper


%---------------------------------------------------------------------------------
% Additional settings
%---------------------------------------------------------------------------------

%---------------------------------------------------------------------------------
% Define your margins
\usepackage{geometry} % Necessary package for defining margins

\geometry{
	top=2cm, % Defines top margin
	bottom=2cm, % Defines bottom margin
	left=2.2cm, % Defines left margin
	right=2.2cm, % Defines right margin
	includehead, % Includes space for a header
	%includefoot, % Includes space for a footer
	%showframe, % Uncomment if you want to show how it looks on the page
}

\setlength{\parindent}{15pt} % Adjust to set you indent globally

%---------------------------------------------------------------------------------
% Define your spacing
\usepackage{setspace} % Required for spacing
% Two options:
\linespread{1.5}
%\onehalfspacing % one-half-spacing linespread

%----------------------------------------------------------------------------------------
% Define your fonts
\usepackage[T1]{fontenc} % Output font encoding for international characters
\usepackage[utf8]{inputenc} % Required for inputting international characters

\usepackage{XCharter} % Use the XCharter font


%---------------------------------------------------------------------------------
% Define your headers and footers

\usepackage{fancyhdr} % Package is needed to define header and footer
\pagestyle{fancy} % Allows you to customize the headers and footers

%\renewcommand{\sectionmark}[1]{\markboth{#1}{}} % Removes the section number from the header when \leftmark is used

% Headers
\lhead{} % Define left header
\chead{\textit{}} % Define center header - e.g. add your paper title
\rhead{} % Define right header

% Footers
\lfoot{} % Define left footer
\cfoot{\footnotesize \thepage} % Define center footer
\rfoot{ } % Define right footer

%---------------------------------------------------------------------------------
%	Add information on bibliography
\usepackage{natbib} % Use natbib for citing
\usepackage{har2nat} % Allows to use harvard package with natbib https://mirror.reismil.ch/CTAN/macros/latex/contrib/har2nat/har2nat.pdf

% For citing with natbib, you may want to use this reference sheet:
% http://merkel.texture.rocks/Latex/natbib.php

%---------------------------------------------------------------------------------
% Add field for signature (Reference: https://tex.stackexchange.com/questions/35942/how-to-create-a-signature-date-page)
\newcommand{\signature}[2][5cm]{%
  \begin{tabular}{@{}p{#1}@{}}
    #2 \\[2\normalbaselineskip] \hrule \\[0pt]
    {\small \textit{Signature}} \\[2\normalbaselineskip] \hrule \\[0pt]
    {\small \textit{Place, Date}}
  \end{tabular}
}

\renewcommand{\abstractname}{摘要}
%---------------------------------------------------------------------------------
%	General information
%---------------------------------------------------------------------------------
\title{Project Report} % Adds your title
\author{
Qi Gao \big(qxg150130@utdallas.edu\big)\\
Pancham Mamania \big(pxm172730@utdallas.edu\big)
 % Add your first and last name
    %\thanks{} % Adds a footnote to your title
    %\institution{YOUR INSTITUTION} % Adds your institution
  }

\date{12/01/2019} % Adds the current date to your “cover” page; leave empty if you do not want to add a date


%---------------------------------------------------------------------------------
%	Define what’s in your document
%---------------------------------------------------------------------------------

\begin{document}
\lstset{numbers=left,
numberstyle=\tiny,
keywordstyle=\color{blue!70},
commentstyle=\color{red!50!green!50!blue!50},
%frame=shadowbox,
rulesepcolor=\color{red!20!green!20!blue!20},
breaklines=true,
extendedchars=true
}

% If you want a cover page, uncomment "%---------------------------------------------------------------------------------
% Cover page
%---------------------------------------------------------------------------------

% Here are more templates for other cover pages: https://www.latextemplates.com/cat/title-pages

% This example is based on this cover page example: https://www.latextemplates.com/template/academic-title-page

\begin{titlepage} % Starts new environment where the page number is not displayed and the count starts at 1 for the next page

%------------------------------------------------
%	Institutional information
%------------------------------------------------
	
\begin{minipage}{0.4\textwidth} % Begins new environment (like a text box)
    \begin{flushleft} % Sets environment on the left side of the paper
    \large
    University of XX\\ % Add your institution
    Chair of Political Science IV\\ % Add the chair
    Fall 2018\\ % Add term
    COURSE TITLE\\ % Add course title
    Supervisor: NAME % Add instructor/supervisor name 
    \end{flushleft}
\end{minipage}
	
\vspace*{2in} % Adds some space in-between
	
\center % Centre everything on the page

%------------------------------------------------
%	Main part
%------------------------------------------------
	
{\huge\bfseries TITLE OF YOUR PAPER}\\[0.4cm] % Add your paper title 
{\large\today}\\[0.4cm] % Add date (current day)
FIRSTNAME LASTNAME % Add your name
	
\vfill % Adds additional space

%------------------------------------------------
%	General information about the author
%------------------------------------------------

\vfill % Adds additional space

Your contact info \\ % Add your contact info
Your Program \\ % Add info about your program
Semester you are enrolled \\ % Add info about your semester

\vfill % Adds additional space

%------------------------------------------------
%	Word count
%------------------------------------------------

\vfill % Adds additional space
	
Word count: XXXX % To indicate the word count
% How to check words in a LaTeX document: https://www.overleaf.com/help/85-is-there-a-way-to-run-a-word-count-that-doesnt-include-latex-commands
	

	
\end{titlepage}" and uncomment "\begin{comment}" and "\end{comment}" to comment the following lines
%%---------------------------------------------------------------------------------
% Cover page
%---------------------------------------------------------------------------------

% Here are more templates for other cover pages: https://www.latextemplates.com/cat/title-pages

% This example is based on this cover page example: https://www.latextemplates.com/template/academic-title-page

\begin{titlepage} % Starts new environment where the page number is not displayed and the count starts at 1 for the next page

%------------------------------------------------
%	Institutional information
%------------------------------------------------
	
\begin{minipage}{0.4\textwidth} % Begins new environment (like a text box)
    \begin{flushleft} % Sets environment on the left side of the paper
    \large
    University of XX\\ % Add your institution
    Chair of Political Science IV\\ % Add the chair
    Fall 2018\\ % Add term
    COURSE TITLE\\ % Add course title
    Supervisor: NAME % Add instructor/supervisor name 
    \end{flushleft}
\end{minipage}
	
\vspace*{2in} % Adds some space in-between
	
\center % Centre everything on the page

%------------------------------------------------
%	Main part
%------------------------------------------------
	
{\huge\bfseries TITLE OF YOUR PAPER}\\[0.4cm] % Add your paper title 
{\large\today}\\[0.4cm] % Add date (current day)
FIRSTNAME LASTNAME % Add your name
	
\vfill % Adds additional space

%------------------------------------------------
%	General information about the author
%------------------------------------------------

\vfill % Adds additional space

Your contact info \\ % Add your contact info
Your Program \\ % Add info about your program
Semester you are enrolled \\ % Add info about your semester

\vfill % Adds additional space

%------------------------------------------------
%	Word count
%------------------------------------------------

\vfill % Adds additional space
	
Word count: XXXX % To indicate the word count
% How to check words in a LaTeX document: https://www.overleaf.com/help/85-is-there-a-way-to-run-a-word-count-that-doesnt-include-latex-commands
	

	
\end{titlepage}

%\begin{comment}
\maketitle % Print your title, author name and date; comment if you want a cover page

%\begin{center} % Center text
    %Word count: XXXX
% How to check words in a LaTeX document: https://www.overleaf.com/help/85-is-there-a-way-to-run-a-word-count-that-doesnt-include-latex-commands
%\end{center}
%\end{comment}

%----------------------------------------------------------------------------------------
% Introduction
%----------------------------------------------------------------------------------------
\setcounter{page}{1} % Sets counter of page to 1
\section{Ethereum and Tokens Description} % Add a section title
\subsection{Ethereum and ERC20} % Add a subsection
Ethereum is an innovation that applies some of the technologies and concepts of Bitcoin in computing. Bitcoin is considered a system that maintains a shared global book that securely records all bitcoin bills. Ethereum uses a number of mechanisms similar to Bitcoin (such as blockchain technology and P2P networks) to maintain a shared computing platform that can flexibly and safely run any program the user wants.

\subsection{Project Tokens}  % Add another subsection
\subsubsection{Tierion}
Tierion is building a common data validation platform. Tierion's working principle is simply to create a proof that associates data with transactions on a blockchain. This process is called "anchoring." Anyone with this proof can verify the integrity and time stamp of the data without relying on any trusted authority.

\subsubsection{Aragon}
Aragon's decentralized APP on the Ethereum blockchain allows anyone to create and manage any organization. It implements the basic functions of shareholder roster, token transfer, voting, job appointment, financing, and accounting. It can be defined by modifying the charter. The behavior of organizations on the chain provides opportunities to create and manage decentralized organizations.

\subsubsection{Bitqy}
Bitqy is also called BQ, which is the official cryptocurrency issued by Bitqyck. Bitqyck will use BQ to scan various cryptocurrencies, including bitcoin and full-service data support and management.


%----------------------------------------------------------------------------------------
% Literature review
%----------------------------------------------------------------------------------------

\section{Project Goal}
This project has two parts. The first goal is to filter out all outliers from each tokens and then find out the best distribution for each token. The second part is that we find out the best K that is the number of most active buyers, which bought most tokens or has most number of transactions. And the top K buyers' data gives the best fitted multiple regression model.

\subsection{Question 1}
The tokens we picked is Tierion, Aragon, and Bitqy. Because we cannot find token126, hms and lino in token pirce folder, so we moved down and picked Bitqy as our third token.
First of all, we read data file and then filtered out all outliers which amount of tokens per transcation is larger than total supply.
\lstset{language=R}
\begin{lstlisting}
bitqy <- read_delim('C:/Users/ygaoq/OneDrive/MyDocuments/2019 Spring/Statistics/Project/Blockchain-Tokens-Data-Analytics/networkbitqyTX.txt', delim = " ", col_names = F)
names(bitqy) <- c('fromID', 'toID', 'unixTime', 'tokenAmount')
decimals <- 2
supply <- 1 * 10^10
#filter out all outliers
bitqyFiltered <-filter(bitqy,tokenAmount < decimals * supply)
#figure out how many users indruced those unnormal transcation
bitqy_outliers<- filter(bitqy,tokenAmount >= decimals * supply)
user_outliers <- bitqy_outliers %>% group_by(toID) %>% summarise(n = n()) %>% ungroup
number_users_outliers<-nrow(user_outliers)
\end{lstlisting}

\subsubsection{And a subsubsection}

And here you have a list with dots and dashes:
\begin{itemize}
    \item Dots
    \item[-] or dashes?
\end{itemize}

and an enumerated list:

\begin{enumerate}
    \item With numbers...
    \item Which is also nice \Winkey
\end{enumerate}

%---------------------------------------------------------------------------------
% Theory
%---------------------------------------------------------------------------------

\section{Theory}

If you want to add mathematical equations, this may be done either this way: $a +  b \neq \frac{a}{b}$. You may also add Greek letters like this: $\alpha$.\\

\blindtext % Some blind text

% Including figures
\begin{figure}[htpb!] % Defines figure environment
    \centering % Centers your figure
\includegraphics[scale=0.8]{figure/figure.png} % Includes your figure and defines the size
    \caption{A circle} % For your caption
    \label{fig:my_label} % If you want to label your figure for in-text references
\end{figure}

\blindtext % Some blind text

% Including tables
%   Simple table
\begin{table}[] % Add htpb! to make sure that table is where it should be
    \centering
    \begin{tabular}{c|c}
        Saturday & Sunday \\
        12 & 18
    \end{tabular}
    \caption{Overview of the weekend}% Caption for tables
    \label{tab:weekend} % Reference for in-line referencing
\end{table}

%   Table with table generator
\begin{table}[] % Add htpb! to make sure that table is where it should be
    \centering
\begin{tabular}{@{}lll@{}}
\toprule
  & A   & B \\ \midrule
C & 100 & 2 \\
D & 3   & 5 \\ \bottomrule
\end{tabular}
    \caption{Random numbers} % Caption for tables
    \label{tab:numbers} % Reference for in-line referencing
\end{table}

%----------------------------------------------------------------------------------------
% Research design
%----------------------------------------------------------------------------------------

\section{Research Design}

\blindtext % Some blind text

%----------------------------------------------------------------------------------------
% Analysis
%----------------------------------------------------------------------------------------

\section{Analysis}

\blindtext % Some blind text

%----------------------------------------------------------------------------------------
% Conclusion
%----------------------------------------------------------------------------------------

\section{Conclusion}

\blindtext % Some blind text

%----------------------------------------------------------------------------------------
% Bibliography
%----------------------------------------------------------------------------------------
\newpage % Includes a new page

\pagenumbering{roman} % Changes page numbering to roman page numbers
%\bibliography{literature}

\bibliography{literature.bib} % Add the filename of your bibliography
\bibliographystyle{apsr} % Defines your bibliography style

% For citing, please see this sheet: http://merkel.texture.rocks/Latex/natbib.php

%----------------------------------------------------------------------------------------
% Appendix
%----------------------------------------------------------------------------------------
\newpage % Includes a new page
\section*{Appendix} % Stars disable section numbers
% \appendix % Uncomment if you want to add an "automatic" appendix
\pagenumbering{Roman} % Changes page numbering to Roman page numbers

\blindtext % Adds some blind text

%----------------------------------------------------------------------------------------
% Declaration
%----------------------------------------------------------------------------------------
\newpage % Includes a page break
\thispagestyle{empty} % Leaves the page style empty (no page number, no header, no footer)
\section*{Statutory Declaration} % Stars disable section numbers

\begin{otherlanguage}{german}
Hiermit versichere ich, dass diese Arbeit von mir pers\"{o}nlich verfasst ist und dass ich keinerlei fremde Hilfe in Anspruch genommen habe. Ebenso versichere ich, dass diese Arbeit oder Teile daraus weder von mir selbst noch von anderen als Leistungsnachweise andernorts eingereicht wurden. W\"{o}rtliche oder sinngem\"{a}{\ss}e \"{U}bernahmen aus anderen Schriften und Ver\"{o}ffentlichungen in gedruckter oder elektronischer Form sind gekennzeichnet. S\"{a}mtliche Sekund\"{a}rliteratur und sonstige Quellen sind nachgewiesen und in der Bibliographie aufgef\"{u}hrt. Das Gleiche gilt f\"{u}r graphische Darstellungen und Bilder sowie f\"{u}r alle Internet-Quellen. Ich bin ferner damit einverstanden, dass meine Arbeit zum Zwecke eines Plagiatsabgleichs in elektronischer Form anonymisiert versendet und gespeichert werden kann. Mir ist bekannt, dass von der Korrektur der Arbeit abgesehen und die Pr\"{u}fungsleistung mit nicht ausreichend bewertet werden kann, wenn die Erkl\"{a}rung nicht erteilt wird.
\end{otherlanguage}

\vspace*{1in} % Adds extra space between two paragraphs

\noindent I hereby declare that the paper presented is my own work and that I have not called upon the help of a third party. In addition, I affirm that neither I nor anybody else has submitted this paper or parts of it to obtain credits elsewhere before. I have clearly marked and acknowledged all quotations or references that have been taken from the works of others. All secondary literature and other sources are marked and listed in the bibliography. The same applies to all charts, diagrams and illustrations as well as to all Internet resources. Moreover, I consent to my paper being electronically stored and sent anonymously in order to be checked for plagiarism. I am aware that the paper cannot be evaluated and may be graded ``failed'' (``nicht ausreichend'') if the declaration is not made.\\

%\vspace*{1in} % Adds extra space

% Add field for signature, date, and place
\hfill \signature{}


%---------------------------------------------------------------------------------

\end{document}
